\chapter{Conclusion}
\label{chap:conclusion}

This thesis asked whether spike-sorting quality metrics can be predicted for real paired electrophysiology recordings by transferring supervision from larger hybrid SpikeForest datasets. The answer is qualified but clear.

The benchmark shows that transfer is difficult for principled reasons. Hybrid and paired data are strongly shifted, and naive source-only transfer fails. Once that shift is respected rather than ignored, context-aware transfer models produce a credible positive benchmark for \fpos{}, but the main robust story is narrower than one final model wins. On the twenty-seed recording-disjoint ladder, the strongest overall \fpos{} score in the retained core set is the paired-only raw model at 0.4209 \rsq{}, while among the structured transfer stacks the pre-waveform unlabeled-structure model is the strongest robust transfer reference at 0.4068 mean \rsq{}. Under family-held-out LOFO evaluation, the waveform-augmented structure model becomes slightly stronger overall and wins three of the five held-out families, but again not universally.

The same cannot yet be said for \fmiss{}. Context-aware models repair some of the damage done by pure source transfer, but they do not cross into a stable positive regime. In the retained core ladder, the reduced-latent context model reaches only -0.0128 mean \rsq{} across the twenty-seed sweep, and under LOFO the least-bad aggregate recipe falls to -0.3673. The thesis therefore concludes that the current feature space and transfer setup are sufficient for a strong \fpos{} benchmark contribution, but not yet for a robust \fmiss{} deployment story.

\section{Final Resolution of the Hypotheses}

\begin{itemize}
    \item \textbf{H1 supported.} The hybrid and paired domains are strongly shifted, and this shift is measurable through near-perfect domain separability.
    \item \textbf{H2 supported.} Leakage-safe recording context adds meaningful predictive value and becomes the dominant explanatory block in the canonical \fpos{} model.
    \item \textbf{H3 supported.} \fpos{} and \fmiss{} are structurally different prediction targets and should not be forced into one shared modelling story.
    \item \textbf{H4 supported.} The preferred \fpos{} recipe depends on study family and on whether the evaluation objective is recording-disjoint robustness or unseen-family robustness.
\end{itemize}

\section{Main Takeaway}

The main takeaway is that leakage-safe spike-sorting quality prediction under hybrid-to-paired transfer is feasible for \fpos{}, not yet robust for \fmiss{}, and strongly conditioned by target-domain heterogeneity. In practical terms, the project is strongest as a benchmark and analysis contribution: it shows where quality prediction works, where it remains fragile, and why the answer depends on both the biological regime and the evaluation regime.
